\documentclass[12pt,a4paper]{article}
\usepackage[utf8]{inputenc}
\usepackage[spanish,es-noquoting]{babel}
\usepackage[T1]{fontenc}
\usepackage{lmodern}
\usepackage{geometry}
\geometry{top=2cm, bottom=2cm, left=2.5cm, right=2.5cm}
\usepackage{xcolor}
\usepackage{listings}
\usepackage{enumitem}
\usepackage{fancyhdr}
\usepackage{hyperref}
\usepackage{tikz}
\usetikzlibrary{arrows.meta, positioning, calc}

\definecolor{codegreen}{rgb}{0,0.6,0}
\definecolor{backcolour}{rgb}{0.95,0.95,0.92}

\lstset{
    language=C,
    backgroundcolor=\color{backcolour},
    commentstyle=\color{codegreen},
    keywordstyle=\color{blue},
    basicstyle=\ttfamily\small,
    breaklines=true,
    frame=single,
    numbers=left,
    tabsize=4,
    extendedchars=true,
    showstringspaces=false,
    literate={á}{{\'a}}1 {é}{{\'e}}1 {í}{{\'i}}1 {ó}{{\'o}}1 {ú}{{\'u}}1 {ñ}{{\~n}}1
}

\pagestyle{fancy}
\fancyhf{}
\lhead{Monitorías Sistemas Operativos 2026-I}
\rhead{Capítulo 0: Vector Dinámico}
\cfoot{\thepage}

\title{\textbf{Guía de Aprendizaje: Creación de un Vector Dinámico}}
\author{Malak Sanchez \\ \small Monitorías de Sistemas Operativos}
\date{}

\begin{document}

\maketitle

\section{Introducción: Stack vs Heap}
Para entender la memoria dinámica, primero debemos diferenciar dónde vive cada variable en nuestro programa.

\begin{itemize}
    \item \textbf{Stack (Pila):} Es memoria automática y rápida. Aquí viven las variables locales (como \texttt{int x}). Su tamaño se define al compilar. Cuando la función termina, esta memoria se libera automáticamente.
    \item \textbf{Heap (Montículo):} Es un gran bloque de memoria libre gestionado por el programador. Aquí vive la memoria dinámica. Es más lenta que el stack, pero su tamaño puede decidirse \textbf{mientras el programa corre}. Tú eres responsable de pedirla y liberarla.
\end{itemize}

\begin{center}
\begin{tikzpicture}[node distance=2cm, font=\small]
    % Stack
    \draw[thick, fill=blue!10] (0,0) rectangle (3,4);
    \node[above] at (1.5, 4) {\textbf{STACK}};
    \draw[fill=white] (0.2, 3.2) rectangle (2.8, 3.7) node[midway] {\texttt{int *array}};
    \draw[fill=white] (0.2, 2.5) rectangle (2.8, 3.0) node[midway] {\texttt{int n}};
    
    % Heap
    \draw[thick, fill=green!10] (6,0) rectangle (10,4);
    \node[above] at (8, 4) {\textbf{HEAP}};
    \foreach \i in {0,1,2,3} {
        \draw[fill=white] (6.5, 3.5 - \i*0.8) rectangle (9.5, 3.0 - \i*0.8) node[midway] {Dato \i};
    }
    
    % Arrow
    \draw[-{Stealth}, thick, red] (2.8, 3.45) -- (6.5, 3.25) node[midway, above, sloped] {Apunta al Heap};
    
    \node[below=0.2cm] at (1.5,0) {Memoria Automática};
    \node[below=0.2cm] at (8,0) {Memoria Dinámica};
\end{tikzpicture}
\end{center}

\section{Creación de un Vector Dinámico}
Un \textbf{vector dinámico} es simplemente un arreglo cuyo tamaño se define en tiempo de ejecución.

\subsection{El proceso en 3 pasos}
\begin{enumerate}
    \item \textbf{Declarar un puntero:} El puntero vivirá en el \textit{Stack}.
    \item \textbf{Pedir memoria:} Usamos \texttt{malloc} para reservar espacio en el \textit{Heap}.
    \item \textbf{Liberar memoria:} Usamos \texttt{free} cuando ya no necesitemos el arreglo.
\end{enumerate}

\section{Código de Ejemplo}
\begin{lstlisting}
#include <stdio.h>
#include <stdlib.h>

int main(){
    int n;
    printf("Enter vector size: ");
    scanf("%d", &n);

    // 1 and 2. Declare pointer and request memory in the Heap
    int *array = (int *) malloc(n * sizeof(int));

    // Check if allocation was successful
    if(array == NULL){
        printf("Error: Could not allocate memory.\n");
        return 1;
    }

    // Use the array as a normal array
    for(int i=0; i < n; i++){
        array[i] = i * 10;
        printf("array[%d] = %d\n", i, array[i]);
    }

    // 3. Free memory when finished
    free(array);
    array = NULL; // Good practice

    return 0;
}
\end{lstlisting}

\section{¿Por qué usar malloc y sizeof?}
La función \texttt{malloc} recibe la cantidad de \textbf{bytes} a reservar. Como no sabemos cuánto pesa un \texttt{int} en todos los sistemas, usamos \texttt{sizeof(int)}. Así, \texttt{n * sizeof(int)} garantiza espacio exacto para $n$ enteros.

\section{Resumen Visual}
Cuando haces \texttt{int *v = malloc(3 * sizeof(int))}:
\begin{enumerate}
    \item En el \textbf{Stack} se crea \texttt{v}.
    \item En el \textbf{Heap} se buscan $3 \times 4 = 12$ bytes contiguos.
    \item \texttt{v} guarda la dirección del primer byte de ese bloque en el Heap.
\end{enumerate}

\end{document}
